\documentclass[11pt]{article}
\usepackage[margin=1in]{geometry}
\usepackage{amsmath,amssymb,amsthm}
\usepackage{enumitem,booktabs}
\usepackage[hyphens]{url}
\usepackage[hidelinks]{hyperref}
\setlist{itemsep=1pt,topsep=3pt}
\emergencystretch=2em
\newtheorem{theorem}{Theorem}[section]
\newtheorem{lemma}[theorem]{Lemma}
\newtheorem{corollary}[theorem]{Corollary}
\newtheorem{proposition}[theorem]{Proposition}
\theoremstyle{definition}
\newtheorem{definition}[theorem]{Definition}
\newtheorem{remark}[theorem]{Remark}
\newcommand{\st}[1]{\textsf{\small[#1]}}
\newcommand{\file}[1]{\path{#1}}
\newcommand{\where}[1]{\par\noindent{\small\emph{In the repo:} #1}\par\smallskip}
\newcommand{\F}{\mathbb{F}_2}
\newcommand{\wt}{\operatorname{wt}}
\newcommand{\codim}{\operatorname{codim}}
\newcommand{\binomV}[1]{\tbinom{V}{#1}}
\DeclareMathOperator{\ANTI}{ANTI}

\title{Erd\H{o}s Problem \#644: progress report\\[3pt]
\large barrier results, an exact classification, and a reduction to one open statement}
\author{erdos644 project}
\date{25 November 2026 (draft)}

\begin{document}
\maketitle

\begin{abstract}
\noindent This is a progress write-up for readers of the repository, not a paper. We do \textbf{not} resolve
the problem: we neither prove $f(k,7)\le(3/4+o(1))k$ nor disprove it. We record what has been proved, at which level of
checking, and the single statement to which the problem now reduces. Every result carries a status tag:
\st{LEAN} machine-checked; \st{VERIFIED} hand proof that passed an independent re-derivation; \st{EXACT} exact
(rational or exhaustive) computation; \st{HAND} hand proof not yet independently checked. Tags are copied from
\file{notes/LEDGER.md}, which is the authority; three short consequences derived here from verified results are
tagged as such.
\end{abstract}

\section{Where things stand}\label{sec:summary}

\paragraph{The problem.} For a finite family $H$ of sets, $\tau(H)$ is the least size of a set meeting every member.
$H$ has \emph{property $(7,2)$} if every subfamily of \emph{at most} seven members has a transversal of at most two
points. (With ``exactly seven distinct members'' families with fewer than seven members are vacuously $(7,2)$; the
two conventions agree for families with at least seven members, and the asymptotic question is the same.) Let
$f(k,7)$ be the largest $\tau$ of a $k$-uniform $(7,2)$-family. Erd\H{o}s Problem \#644~\cite{EP} asks whether
$f(k,7)=(3/4+o(1))k$.

\paragraph{Known bounds (unchanged by this work).} For property $(p,2)$ with $p=3,4,5,6$ the exact values are
$2k,\lfloor 3k/2\rfloor,\lfloor 5k/4\rfloor,k$~\cite{EFKT}. For $p=7$, Fon-Der-Flaass, Kostochka and
Woodall~\cite{FKW} proved
\[
3m+1\le f(4m,7)\ (m\ge10),\qquad f(k,7)\le\lceil 7k/8\rceil ,
\]
so $\tfrac34k+O(1)\le f(k,7)\le\tfrac78k+1$. Their lower-bound family is a parity family; see \S\ref{sec:linear}.

\paragraph{What we prove.} Write $T=\lceil(3/4+\varepsilon)k\rceil$, $n=|\bigcup H|=(7/4+\delta)k$ and $N=n/k$.
\begin{enumerate}[leftmargin=*]
\item \emph{Reformulation} (\S\ref{sec:setting}). Sets of size at most $k$ give the same $f$; the problem is
Statement (G) below. A seven with no 2-point transversal is produced by a ``Fano colouring'' \st{VERIFIED, LEAN}.
\item \emph{Exact classification} (\S\ref{sec:class}). If $n<11k/6$, every seven members with no 2-point
transversal are of Fano type. The constant $11/6$ is exact for the relaxation \st{EXACT}.
\item \emph{Barrier results} (\S\ref{sec:linear}, \S\ref{sec:coset}). The pencil-sum identity explains FKW in three
lines. Conversely, for covering coset families in the Fano range the parity obstruction can certify $(7,2)$ only when
$\tau\le 3k/4+3/2$ (Theorem~D, ``$c^*=1$''). This is why parity has never beaten $3k/4+O(1)$ \st{VERIFIED}.
\item \emph{Structured families lose} (\S\ref{sec:rich}--\S\ref{sec:coset}). Rich families (Theorem RA), families
with a near-complete block (Theorem NC) and, in the Fano range, coset families with few classes (Theorems D and
$9'$) are not $(7,2)$ once $\tau>(3/4+\varepsilon)k$ and $k$ is large \st{VERIFIED}. More coset ranges are covered by HAND results.
\item \emph{Small values} (\S\ref{sec:small}). $f(6,7)=5$, $f(7,7)=6$, $f(8,7)\in\{6,7\}$ \st{EXACT}.
\end{enumerate}

\paragraph{The four-region map.} Pad to $k$-uniform and let $H$ be a putative counterexample with
$\tau>(3/4+\varepsilon)k$. By Theorem~\ref{thm:leanA} \st{LEAN} it has $n\ge(7/4+3\varepsilon)k-9$.
\begin{center}
\begin{tabular}{@{}lll@{}}
\toprule
Region & Content & Status\\
\midrule
I.\ Rich: $\tau^*_b(H)\le 7/3-\eta$ & Theorem RA & closed \st{VERIFIED}\\
II.\ A near-complete block & Theorem NC & closed \st{VERIFIED}\\
III.\ Linear layers (cosets) & D, $9'$; B1, S0, S1 & partly closed \st{VERIFIED/HAND}\\
IV.\ Thin, no block, not linear & Statement (S) & \textbf{open}\\
\bottomrule
\end{tabular}
\end{center}
Region IV is the whole difficulty. Proposition~\ref{prop:S} shows that (G) is \emph{equivalent} to:

\begin{quote}
\textbf{(S)} For some fixed $\eta\in(0,4/3]$ and every $\varepsilon>0$ there is $k_0$ such that for $k\ge k_0$:
every $k$-uniform family $H$ with $\tau(H)\ge(3/4+\varepsilon)k+1$ and $\tau^*_b(H)>7/3-\eta$, which contains no
$(W,d)$-complete block with $|W|\ge(7/4-\varepsilon)k+2d+3$, has seven members with no transversal of size at
most~$2$.
\end{quote}
Here $H$ is \emph{$(W,d)$-complete} if every $(k+d)$-subset of $W$ contains a member. For $n<11k/6$ the seven may be taken of Fano type (\S\ref{sec:class}); for $n\ge 11k/6$ nothing general is known.
Computational evidence (\S\ref{sec:open}) suggests that (S) is a supersaturation statement for sparse families
rather than a structure statement.

\section{Setting}\label{sec:setting}

\begin{lemma}[padding]\label{lem:pad}
Let $f_{\le}(k,7)$ be the largest $\tau$ of a $(7,2)$-family of nonempty sets of size at most $k$. Then
$f_{\le}(k,7)=f(k,7)$. \st{VERIFIED, LEAN}
\end{lemma}
\begin{proof}
Add $k-|B|$ fresh private points to each member $B$. A 2-point transversal of $\le7$ original members still meets the
padded members, so $(7,2)$ survives. Conversely, in a transversal of the padded family replace each private point of $B'$
by a point of $B$: this gives a transversal of the original family, no larger. So $\tau$ is unchanged.
\end{proof}
\where{\file{notes/MINIMAL_FRAME.md}; \file{notes/verify/minimal_frame.md}; \file{Lean644/MinimalFrame.lean}
(\texttt{padding\_iff}).}

\noindent\textbf{Statement (G).} \emph{For every $\varepsilon>0$ and $k\ge k_0(\varepsilon)$, every finite family of
nonempty sets of size at most $k$ with $\tau>\lceil(3/4+\varepsilon)k\rceil$ has at most seven members with no
transversal of size at most~$2$.} By Lemma~\ref{lem:pad}, (G) is equivalent to $f(k,7)\le(3/4+o(1))k$
\st{VERIFIED}. There are no anchors, no confinement and no bound on $n$.

\paragraph{Fano colourings.} Let $\Pi$ be the Fano plane. A \emph{colouring} is a map $\chi:V\to\Pi$ from
$V=\bigcup H$ to the points. For a line $\ell$ put $L_\ell=\chi^{-1}(\ell)$ (points whose colour lies on $\ell$) and
$R_\ell=V\setminus L_\ell$. A \emph{Fano-type seven} is a choice of members $B_\ell\subseteq R_\ell$, one per line
(repeats allowed), for some colouring. Equivalently, $x\notin B_\ell$ whenever $\chi(x)\in\ell$.

\begin{lemma}[colouring criterion]\label{lem:col}
If no $L_\ell$ is a transversal of $H$, then $H$ has a Fano-type seven, and it has no transversal of at most two
points. So $H$ is not $(7,2)$. \st{VERIFIED, LEAN}
\end{lemma}
\begin{proof}
Pick $B_\ell\subseteq R_\ell$. For $x,y\in V$ with colours $p,q$, some line $\ell$ contains $p$ and $q$; then
$x,y\in L_\ell$ and $B_\ell$ misses both.
\end{proof}

\begin{theorem}[weighted static criterion, Theorem 3H]\label{thm:3H}
Let $y:V\to\mathbb R_{\ge0}$ satisfy $y(M)\ge1$ for every transversal $M$. Let $A\subseteq V$, $a'=\max_{V\setminus A}y$,
$\Lambda=y(V\setminus A)$, and suppose there are $\pi:A\to\Pi$ and $c\in\mathbb R^{\Pi}_{\ge0}$ with
$\sum_pc_p=\Lambda$ and, for every line $\ell$,
$y(\pi^{-1}(\ell))+\sum_{p\in\ell}c_p+3a'<1$. Then $H$ is not $(7,2)$. \st{VERIFIED}
\end{theorem}
\begin{proof}
Colour $A$ by $\pi$. Feed the points of $V\setminus A$ in any order into class $p_1$ until its load reaches
$c_{p_1}$, then into $p_2$, and so on. Every class ends with load $\le c_p+a'$. So $y(L_\ell)<1$, $L_\ell$ is not a
transversal, and Lemma~\ref{lem:col} applies.
\end{proof}
With $A=\emptyset$ and $c_p=\Lambda/7$ this gives \emph{Theorem 3}: $y(V)+7\max y<7/3$ implies not $(7,2)$. With
$y\equiv1/\tau$ every $(7,2)$-family on $n$ points has $\tau\le3(n+7)/7$ \st{VERIFIED}. The fractional parameter is
$\tau^*_b(H)=\min\{y(V):y\ge0,\ y(M)\ge1$ for every transversal $M\}$. $H$ is \emph{rich} if
$\tau^*_b\le7/3-\eta$ and \emph{thin} otherwise.
\where{\file{notes/PROGRAMME_2.md}; \file{notes/verify/programme2_task0.md} \S1--2;
\file{Lean644/RichFamilies.lean} (\texttt{colouring\_criterion}).}

\begin{theorem}[one-anchor template]\label{thm:leanA}
A $k$-uniform $(7,2)$-family with all members in an $n$-set has $\tau\le\lceil(n-k)/3\rceil+2\lceil k/4\rceil$. In
particular a counterexample to (G) has $n\ge(7/4+3\varepsilon)k-9$. \st{LEAN}
\end{theorem}
\emph{Sketch.} Fix a member $A$; split $A$ into quarters and $V\setminus A$ into thirds; six requests (one third plus two
quarters, along the three perfect matchings of the quarters) together with $A$ force a seven with no 2-point
transversal. Tight for $K^{(4m)}_{7m-1}$ and for FKW.
\where{\file{Lean644/TheoremA.lean}.}

\section{The exact classification}\label{sec:class}

For seven sets $B_1,\dots,B_7$ and a point $x$ put $I_x=\{i:x\notin B_i\}$ (the \emph{co-type}). A pair $\{x,y\}$
($x=y$ allowed) is a transversal iff $I_x\cap I_y=\emptyset$. So the seven have no 2-point transversal iff the
co-types form an \emph{intersecting} family. Fano type means: for some Fano plane $P$ on the index set $[7]$, every
$I_x$ contains a line of $P$ (index the sets by the lines of $\Pi$; the pencils of $\Pi$ are the lines of $P$).

\begin{theorem}\label{thm:class}
Let $B_1,\dots,B_7$ (repeats allowed) each have at least $k$ points, $|\bigcup B_i|<11k/6$, and no transversal of size
at most $2$. Then they form a Fano-type seven. The constant $11/6$ is exact for the fractional relaxation below.
\st{VERIFIED, EXACT; computer-assisted}
\end{theorem}
\begin{proof}
For a family $M$ of subsets of $[7]$ let
$\mathrm{LP}(M)=\min\{\sum_{F\in M}m_F:\ m\ge0,\ \sum_{F\in M,\,F\not\ni i}m_F\ge1\ (i\in[7])\}$.
The co-type support $S=\{I_x\}$ is intersecting, so it lies in a maximal intersecting family $M$. The masses
$m_F=|\{x\in\bigcup B_i:I_x=F\}|/k$ are feasible for $\mathrm{LP}(M)$, since $\sum_{F\not\ni i}m_F=|B_i|/k\ge1$, with value
$|\bigcup B_i|/k<11/6$. For a Fano plane $P$ the up-set $U_P$ of its lines is maximal intersecting: two lines meet,
and since the Fano plane is not 2-colourable, exactly one of $G$ and $[7]\setminus G$ contains a line. So it suffices
that $\mathrm{LP}(M)\ge11/6$ for every maximal intersecting $M$ other than the 30 families $U_P$.

\emph{Computation.} All self-dual up-sets of $2^{[7]}$ were enumerated: $1\,422\,564$ of them, matching the known
number of maximal intersecting families on 7 points (OEIS A001206~\cite{OEIS}), in 716 classes under $S_7$ (the
orbit--stabiliser sum equals the total). Each class LP was solved and re-certified with a rational primal and a
rational dual, both checked exactly, with equal objectives. Results: the dictator class is infeasible; the Fano class
has value $7/4$; every other class has value $\ge11/6$, with exactly four classes at $11/6$ (dual
$(\tfrac16,\tfrac16,\tfrac16,\tfrac13,\tfrac13,\tfrac13,\tfrac13)$); the next value is $15/8$. An independent
counterexample-guided SAT search (43 exact dual certificates, then UNSAT) confirms the bound.
\end{proof}
At value $11/6$ the optimal supports are still Fano; adding small mass on a line-free co-type of one of the four
extremal classes gives non-Fano patterns with $N$ arbitrarily close to $11/6$. Method A needs no solver trust beyond
the enumeration (whose count matches OEIS); method B rests on an unchecked SAT UNSAT.
\where{\file{notes/verify/programme2_task0.md} \S3; code \file{compute/programme2/verify/maxint7_enum.js},
\file{maxint7_lp.py}, \file{maxint7_sat.py}; log \file{compute/programme2/verify/out/maxint7_lp.log}.}

\section{Linear structure, and why FKW stops at \texorpdfstring{$3k/4+O(1)$}{3k/4+O(1)}}\label{sec:linear}

Identify subsets of $V$ with vectors of $\F^V$: $+$ is symmetric difference and $\wt$ is size.

\begin{lemma}[pencil-sum identity]\label{lem:pencil}
For a Fano-type seven $B_\ell\subseteq R_\ell$, let $s_p=\sum_{\ell\ni p}B_\ell$ and $j_x=|\{\ell:x\in B_\ell\}|$. Then
\[
\sum_{p\in\Pi}\wt(s_p)=\sum_{x}j_x(4-j_x)\ \le\ 3\Big(4\big|\textstyle\bigcup_\ell B_\ell\big|-\sum_\ell|B_\ell|\Big).
\]
If the right side is $0$, all $s_p=0$ and $\{\emptyset\}\cup\{B_\ell\}$ is a linear subspace. \st{VERIFIED}
\end{lemma}
\begin{proof}
Let $J_x=\{\ell:x\in B_\ell\}$ and $c=\chi(x)$. Then $J_x$ avoids the three lines through $c$, so $j_x\le4$, and $x\in
s_p$ iff an odd number of lines of $J_x$ pass through $p$. Count such $p$: one line gives 3; two lines give 4; three
lines avoiding $c$ are not concurrent (the pencil of any $q\ne c$ contains the line $qc$), so they form a triangle with 3 odd
points; the four lines avoiding $c$ cover every point $\ne c$ exactly twice, giving 0. So the count is $j(4-j)$. Then
$j(4-j)\le3(4-j)$ for $1\le j\le4$; sum over $x\in\bigcup B_\ell$. If all $s_p=0$, index lines by
$(\F^3)^*\setminus0$: the pencils are the triples $\{u,v,u+v\}$, so $u\mapsto B_u$ (with $B_0=\emptyset$) is linear.
\end{proof}

A \emph{coset family} is $H\subseteq a+C$ with $C\le\F^V$ linear and $a\notin C$. Let $\rho_{\rm par}$ be the least
weight in $a+C$ that is $\equiv k\pmod2$.

\begin{corollary}[coset obstruction]\label{cor:7}
A $k$-uniform coset family on $n$ points with $7\rho_{\rm par}>3(4n-7k)$ has no Fano-type seven. If also $n<11k/6$,
it is $(7,2)$. \st{VERIFIED}
\end{corollary}
\begin{proof}
Each $s_p$ is a sum of three members, so $s_p\in a+C$ and $\wt(s_p)\equiv3k\equiv k$. Hence $7\rho_{\rm par}\le
\sum\wt(s_p)\le3(4n-7k)$. The second claim is Theorem~\ref{thm:class}.
\end{proof}

\paragraph{FKW in three lines.} Let $|S|=7m+1$, $X\subseteq S$ with $|X|=4m$, $Y=S\setminus X$, and let $H$ be the
$4m$-subsets $B$ with $|B\cap X|$ odd (so $|B\cap Y|$ is odd too). (1) $H$ lies in the coset ``odd on $X$, odd on
$Y$'', where $\rho=2$. (2) $4n-7k=4$, and $7\cdot2=14>12$, so there is no Fano-type seven. (3) $n/k=7/4+1/(4m)<11/6$ for
$m\ge4$, so by Theorem~\ref{thm:class} $H$ is $(7,2)$. Finally $\tau(H)=3m+1$: a $4m$-set even on $X$ has complement
of size $3m+1$ meeting every member, and every $(4m+1)$-set meets both $X$ and $Y$ and so contains an odd member (drop a
$Y$-point or an $X$-point). This recovers FKW's theorem for $m\ge4$, the range of their closing remark
\st{VERIFIED, as a consequence}.

\paragraph{No perfect pencil.} In a coset family any three members have $x_1+x_2+x_3\in a+C$, so
$\wt(x_1+x_2+x_3)\ge\rho_{\rm par}\ge1$. Exact pencils, and the zero-slack sevens of Lemma~\ref{lem:pencil}, never
occur. This is exactly what lets parity families survive at $4n-7k=O(1)$. Corollary~\ref{cor:barrier} shows it cannot do
more.

\begin{remark}[seam example]\label{rem:seam}
Let $4\mid k$, $W=W_1\sqcup W_2$ with $|W|=7k/4$ and $0<|W_1|,|W_2|\le k$, and let $E_k$ be the $k$-subsets odd on
$W_1$. Corollary~\ref{cor:7} with $4n-7k=0$ shows $E_k$ is $(7,2)$, and $\tau(E_k)=3k/4$, as for
$K^{(k)}_{7k/4-1}$. It is both a coset family and $(W,1)$-complete, so the two structures meet at $\varepsilon=0$
\st{VERIFIED}.
\end{remark}

\begin{lemma}[quotient certificate, Lemma Q]
Let $\sigma:\F^V\to Q$ be linear, $\mathcal T\subseteq Q$, $\sigma(B)\in\mathcal T$ for every member, and
$W_\sigma=\min\{\wt(y):\sigma(y)\in\mathcal T+\mathcal T+\mathcal T\}$. If $7W_\sigma>3(4n-7k)$ there is no Fano-type
seven. \st{HAND}
\end{lemma}
\emph{Sketch.} As in Corollary~\ref{cor:7}, with each pencil sum in $\mathcal T+\mathcal T+\mathcal T$. Corollary~\ref{cor:7}
is the case $\sigma=$ quotient by $C$, $\mathcal T=\{\sigma(a)\}$. A Lemma-Q counterexample needs covering together with a gap in the pairwise
intersection spectrum just below $(k-r)/2$; this is problem BAND (\S\ref{sec:open}).
\where{\file{notes/PROGRAMME_2.md}; \file{notes/verify/programme2_task0.md} \S1; \file{notes/verify/theorem_NC.md}
\S4 ($E_k$); \file{notes/PROGRAMME3_CX.md} \S1--2 (Lemma Q, BAND).}

\section{Rich families}\label{sec:rich}

\begin{lemma}[Lemma A0]\label{lem:A0}
If $H$ is a nonempty family of sets of size at most $k$ and every three members (repeats allowed) share a point, then
$\tau(H)\le(k+1)/2$. \st{VERIFIED, LEAN}
\end{lemma}
\begin{proof}
Fix $B_1$. For any $B_2$, $B_1\cap B_2$ meets every member, so $|B_1\cap B_2|\ge\tau$. Hence every
$(|B_1|-\tau+1)$-subset of $B_1$ meets every $B_1\cap B_2$, i.e.\ is a transversal. So $\tau\le k-\tau+1$.
\end{proof}
This is folklore-level: it is Nagy--Patk\'os~\cite{NP}, Prop.~2.2 at $r=3$, followed by the ``intersection minus
points is a cover'' step of Bishnoi--Das--Morris--Szab\'o~\cite{BDMS}, Thm.~1.7. It is tight for
$K^{(k)}_{\lceil3k/2\rceil-1}$. We found no source stating it for non-partite families.

\begin{theorem}[Theorem RA]\label{thm:RA}
Let $\eta>0$ and let $H$ be a family of nonempty sets of size at most $k$ with $\tau^*_b(H)\le7/3-\eta$ and
$\tau(H)>(k+1)/2+49/(3\eta)$. Then $H$ is not $(7,2)$. \st{VERIFIED}
\end{theorem}
\begin{proof}
Take an optimal $y$, put $\theta=\eta/7$, $A=\{y\ge\theta\}$, $\Lambda=y(V\setminus A)$, $a'=\max_{V\setminus A}y<\theta$.
Then $|A|\le y(V)/\theta<49/(3\eta)$ and $y(V)+7a'<7/3$. Also $\tau^*_b\ge1$ forces $\eta\le4/3$, so $a'<1/3$. Label
$\Pi$ so that the lines through a point $0$ are $\{0\}\cup Q_s$ ($s=1,2,3$); each line avoiding $0$ takes one point of
each pair $Q_s$.

\emph{Case $\Lambda+3a'\ge2$.} Apply Theorem~\ref{thm:3H} with $\pi\equiv0$, $c_q=s:=(1-3a')/3-\sigma$ for $q\ne0$ and
$c_0=\Lambda-6s\ge3a'+6\sigma\ge0$, with $\sigma>0$ small. A line avoiding $0$ has mass $3s+3a'=1-3\sigma$. A line
through $0$ has mass $y(A)+c_0+2s+3a'=y(V)-4/3+7a'+4\sigma<1$.

\emph{Case $\Lambda+3a'<2$.} Let $H_A$ be the members disjoint from $A$. Since a transversal of $H_A$ plus $A$ is one of
$H$, $\tau(H_A)>(k+1)/2$, so by Lemma~\ref{lem:A0} some $G_1,G_2,G_3\in H_A$ have empty common intersection. Colour $A$
by $0$; put each $v\notin A$ into a class $X_s$ with $v\notin G_s$; split $X_s$ into two halves of $y$-mass differing by
at most $a'$ and colour them by the two points of $Q_s$. A line through $0$ has $L_\ell=A\cup X_s$, avoided by $G_s$.
A line avoiding $0$ has mass $\le(\Lambda+3a')/2<1$. Lemma~\ref{lem:col} finishes in both cases.
\end{proof}
Only richness and $\tau>k/2+O(1/\eta)$ are used; the threshold $3/4$ plays no role. Richness cannot be weakened to
``$\tau^*_b<7/3$'' at finite $k$: $K^{(6)}_{10}$ is $(7,2)$ with $\tau=5$ and $\tau^*_b=2$.
\where{\file{notes/PROGRAMME2_A.md}; \file{notes/verify/theorem_RA.md}; \file{notes/LIT_A0.md};
\file{Lean644/RichFamilies.lean} (\texttt{lemma\_A0}, \texttt{theorem\_A1}).}

\section{Near-complete blocks}\label{sec:NC}

$H$ is \emph{$(W,d)$-complete} if every $(k+d)$-subset of $W$ contains a member; $w=|W|$.

\begin{lemma}[NC1]\label{lem:NC1}
Let $H$ be $(W,d)$-complete and let $G_1,G_2,G_3\in H$ (repeats allowed) have empty common intersection and
$|W\cap(G_1\cup G_2\cup G_3)|\le2(w-k-d)-3$. Then $H$ is not $(7,2)$. \st{VERIFIED}
\end{lemma}
\begin{proof}
Colour $V\setminus\bigcup G_s$ by $0$; put each $v\in\bigcup G_s$ into some $X_s$ with $v\notin G_s$; split
$X_s\cap W$ into halves of sizes differing by at most $1$ (add $X_s\setminus W$ to either) and colour by $Q_s$. Lines
through $0$ are closed by $G_s$. A line avoiding $0$ has $|L_\ell\cap W|\le(|W\cap\bigcup G_s|+3)/2\le w-k-d$, so
$W\setminus L_\ell$ has $\ge k+d$ points and contains a member. Apply Lemma~\ref{lem:col}.
\end{proof}

\begin{theorem}[Theorem NC]\label{thm:NC}
Let $H$ consist of nonempty sets of size at most $k$, be $(W,d)$-complete with $w\ge k+d+2$, and put
$a=\max(0,2k+2d+3-w)$. If $\tau(H)>(k+1)/2+a$, then $H$ is not $(7,2)$. \st{VERIFIED}
\end{theorem}
\begin{proof}
Take $A\subseteq W$ with $|A|=a$. Then $\tau(H_A)\ge\tau-a>(k+1)/2$, so Lemma~\ref{lem:A0} gives a triple in $H_A$ with
empty intersection. Its $W$-trace lies in $W\setminus A$, of size $w-a\le2(w-k-d)-3$. Apply Lemma~\ref{lem:NC1} to $H$.
\end{proof}

\begin{corollary}\label{cor:NC}
If $\tau\ge\lceil(3/4+\varepsilon)k\rceil+1$ and $H$ is $(W,d)$-complete with $w\ge(7/4-\varepsilon)k+2d+3$, then $H$ is
not $(7,2)$. \st{VERIFIED}
\end{corollary}
The ``$+1$'' matters: with only $\tau>(3/4+\varepsilon)k$ there is a half-point gap for odd $k$ ($K^{(5)}_{10}$,
$K^{(9)}_{17}$; not counterexamples). NC never fires when $n<7k/4+d+3/2$, and it fails to cover a shift-0 layer
whose block excess $d$ is large.
\where{\file{notes/PROGRAMME2_NC.md}; \file{notes/verify/theorem_NC.md}.}

\section{Coset families}\label{sec:coset}

Write $D=a+C=\{x\subseteq V:\sum_{i\in x}h_i=s\}$ with columns $h_i\in\F^M$, $s\ne0$, and $H=D\cap\binomV k$. Let
$r$ be least such that every $(k+r)$-subset of $V$ contains a member, $w=k+r$, and $T=n-w$. Then $\tau(H)=T+1$ (a set
is a transversal iff its complement contains no member). Let $E=\mathrm{span}\{h_i+h_j\}$ and $e=\dim E$.

\begin{theorem}[Theorem D, ``$c^*=1$'']\label{thm:D}
If $n\ge2e+1$ and $e\le\min(w,T)$, then $\rho(a+C)\le r+1$, and $\le r$ when $w$ is even. The same holds for
$\rho_{\rm par}$. In particular this holds when $\dim C\ge k+r>T$. \st{VERIFIED}
\end{theorem}
\begin{proof}
\emph{(A) For $0\ne c\in C$: $\rho\le r+|\wt(c)-w|$.} If $\wt c\ge w$, take a member $x$ inside a $w$-subset of $c$;
then $x+c\in D$ has weight $\wt c-k$. If $\wt c\le w$, take a member $x$ inside a $w$-superset of $c$; then
$|x\cap c|\ge k+\wt c-w$ and $\wt(x+c)\le2w-k-\wt c$.

\emph{(B) If $n\ge2e+1$, there are $e$ disjoint pairs $\{i_t,j_t\}$ with $h_{i_t}+h_{j_t}$ a basis of $E$.} Take a
maximal such family of $t<e$ pairs, spanning $F$. All unmatched columns lie in one coset of $F$ (else add a pair), and
there are $\ge3$ of them. Some matched pair $(i,j)$ has $h_i,h_j$ outside that coset (else $E\subseteq F$). Among
three unmatched positions two, $u,u'$, have $h_u+h_{u'}$ in the span $F'$ of the other pairs (the
$\langle h_i+h_j\rangle$-coefficient is additive). Replacing $(i,j)$ by $(i,u),(j,u')$ gives $t+1$ independent labels, a
contradiction.

\emph{(C) If moreover $e\le w\le n-e$, the sums $\sum Q$ over $w$-sets $Q$ fill $w h_1+E$:} take one endpoint of each
pair plus $w-e$ unmatched positions, and swap endpoints.

\emph{(D)} Here $e\le w=n-T\le n-e$. If $w$ is even, some $w$-set has sum $0$, i.e.\ is a codeword of weight $w$, and
(A) gives $\rho\le r$. If $w$ is odd, some $w$-set $Q$ has sum $h_1$, so $Q\triangle\{1\}$ is a codeword of weight
$w\pm1$, and $\rho\le r+1$. For $\rho_{\rm par}$ append a parity row to the columns: $E$ is unchanged. Finally
$e\le n-\dim C\le n-w=T<w$ and $2e\le2T<n$.
\end{proof}

\begin{corollary}[counting]\label{cor:count}
$|C|\ge|H|\ge\binom nw/\binom{n-k}r$, so $\dim C\ge w$ whenever $\binom nw>2^{w-1}\binom{n-k}r$. This holds for every
$k\le199$, $T\ge\lceil3k/4\rceil$, $n<11k/6$ (exact integer check, $10\,098$ triples), and for all $k$ in that range
from the entropy bound with $g(\omega,t)=Nh(\tfrac{1+\omega}N)-(t+\omega)h(\tfrac\omega{t+\omega})-(1+\omega)\ge0.315$
on $N<11/6$, $t\ge3/4$ (numerical minimum, recomputed independently). Since
$\log_2\binom nw\ge nh(w/n)-\log_2(n+1)$, the entropy bound needs only $0.315k>\log_2(n+1)-1$, true for $k\ge20$; so
the two checks together cover every $k$. At $t=3/4$ the counting reaches $N<1.938$.
\st{VERIFIED}
\end{corollary}

\begin{corollary}[the parity barrier]\label{cor:barrier}
Let $H$ be a covering coset layer as above with $T\ge3k/4$ and $\dim C\ge k+r$ (e.g.\ $n<11k/6$, by
Corollary~\ref{cor:count}). If Corollary~\ref{cor:7} applies to $H$, then $r\le1$ and $4T-3k\le2$, so
$\tau(H)\le3k/4+3/2$. \st{consequence of VERIFIED results}
\end{corollary}
\begin{proof}
$4n-7k=4r+(4T-3k)$. Corollary~\ref{cor:7} needs $7\rho_{\rm par}>12r+3(4T-3k)$, and Theorem~\ref{thm:D} gives
$\rho_{\rm par}\le r+1$. So $3(4T-3k)<7-5r$, which forces $r\le1$ and $4T-3k\le2$.
\end{proof}
So in the Fano range the mechanism behind FKW ($r=1$, $\rho=2$, $4T=3k$) is already at its limit. Coset families with
larger $\tau$ must contain Fano-type sevens; Theorem $9'$ constructs them when the classes are few.

\paragraph{Setting for Theorem $9'$.} Let $Q=\F^V/C$ and $[x]\in Q$ the class of a point. Points of one class form a
\emph{cell}, and any two points of a cell form a weight-2 codeword. Let $\Omega_1,\Omega_2,\Omega_3\in a+C$ be pairwise
disjoint, $w_i=|\Omega_i|$, $W=\sum w_i$, $\Delta=\max w_i-\min w_i$, $V'=V\setminus\bigcup\Omega_i$,
$s=4n-7k-W$, $f=\mathbf 1[s\equiv2\ (4)]$, $u=(s-2f)/4$, and $c$ the number of classes met by $V'$.

\begin{theorem}[Theorem $9'$]\label{thm:9}
Assume (H0) $W\le4n-7k$; (H1) $w_i\equiv k\pmod2$; (H2) $k\ge W+2\Delta+4c+6$; (H3) some class $g^*$ has
$|V'_{g^*}|\ge u+2f+7$; (H4$'$) at most $u$ classes occur exactly once in $V'$. Then $H$ contains seven distinct
members forming a Fano-type seven. \st{VERIFIED}
\end{theorem}
\begin{proof}[Proof sketch]
Fano points $1,\dots,7$, lines $124,235,346,457,156,267,137$. Use the triangles $T_1=\{124,156,267\}$,
$T_2=\{124,137,346\}$, $T_3=\{124,235,457\}$ (contained in the line sets avoiding $q_i=3,5,6$), which cover every line
an odd number of times, and $T^*=\{124,235,346\}$ (avoiding $7$). Remove a set $Z$ of $u$ points (the singleton
classes and filler from $g^*$) and, if $f=1$, a pair $F\subseteq g^*$; call the rest $V''$. Put
$e_\ell=\sum_{i:\ell\in T_i}w_i+2f\mathbf 1[\ell\in T^*]$, $t_\ell=k-e_\ell$ and
$m_p=\tfrac14\sum_\ell t_\ell-\tfrac12\sum_{\ell\ni p}t_\ell$, the solution of $\sum_{p\notin\ell}m_p=t_\ell$. Each
$t_\ell$ is even by (H1), and $\sum_\ell t_\ell=7k-3W-6f\equiv0\pmod4$ by the choice of $f$, so $m_p\in\mathbb Z$. A
direct check gives $4m_p\ge k-W-2\Delta-6f\ge4c$ by (H2), and $\sum_pm_p=|V''|$.

Partition $V''$ into $Y_p$ with $|Y_p|=m_p$ and $[\bigcup_{p\notin\ell}Y_p]=0$ for every line. Since the $t_\ell$ are
even, $(m_p\bmod2)_p$ lies in the kernel of the line-complement matrix, which is the Hamming code $\mathcal H$ spanned
by the lines. For each class $g$ choose $\pi_g\in\mathcal H$: $0$ if $|V''_g|$ is even, a line if it is odd and $g\ne
g^*$, and $\pi_{g^*}=(m\bmod2)+\sum_{g\ne g^*}\pi_g$ (weight $\le7\le|V''_{g^*}|$ by (H3), right parity). Put one
point of $V''_g$ into $Y_p$ for each $p\in\pi_g$, then fill with same-class pairs; the residuals are even and $\ge0$
since $m_p\ge c$. Then $[Y_p]=\sum_g\pi_g(p)\,g$ and $[\bigcup_{p\notin\ell}Y_p]=0$ because $\pi_g\in\mathcal H$.

Set $B_\ell=\bigcup_{p\notin\ell}Y_p\cup\bigcup_{i:\ell\in T_i}\Omega_i\cup(F\text{ if }\ell\in T^*)$. Then
$|B_\ell|=k$ and $[B_\ell]=\alpha$ (odd number of $\Omega$'s), so $B_\ell\in H$. Colour $Y_p$ by $p$, $\Omega_i$ by
$q_i$ and $F\cup Z$ by $7$: then $B_\ell\subseteq R_\ell$, and Lemma~\ref{lem:col} applies. The sets are distinct since
every $Y_p\ne\emptyset$.
\end{proof}

\noindent\textbf{Inputs and limits.} Applying Theorem~\ref{thm:D} on $V$, $V\setminus\Omega_1$ and
$V\setminus(\Omega_1\cup\Omega_2)$ gives three disjoint parity-correct $\Omega_i$ of weight $\le r+1$. Then
$W\le3r+3\le4n-7k$, and the counting holds after the shrinks for $N<1.852$ at $t=3/4$ \st{VERIFIED}. But (H2)--(H4$'$)
hold only when $V'$ has few classes (codimension $O(\log k)$). So Theorems D and $9'$ settle coset families in the Fano
range \emph{with few classes}; the general case is the HAND result B1 below.

\begin{proposition}[the literal Proposition 9 is false]\label{prop:9false}
Without (H1) the conclusion fails: on $V=X\sqcup Y\sqcup\{z_1,z_2\}$ with $|X|=|Y|=5$, $k=5$ and
$a+C=\{S:|S\cap X|,|S\cap Y|\text{ odd}\}$, every member contains exactly one $z_j$, so $\tau\le2$. Yet
$\Omega_i=\{x_i,y_i\}$ are disjoint in $a+C$ with $W=6\le4n-7k=13$. \st{VERIFIED}
\end{proposition}
\where{\file{notes/PROGRAMME2_TASK2.md}; \file{notes/PROGRAMME2_TASK3.md};
\file{notes/verify/theorems_9prime_D.md}.}

\paragraph{Further coset results \st{HAND}.} These are hand proofs with exact per-instance checks of the
constructions; none has passed an independent proof verification.
\begin{itemize}[leftmargin=*]
\item \textbf{B1.} For $N<11/6$ every $T$-covering coset family has a Fano-type seven once
$\codim(C_{\rm ev})-1\le\min\big((k-5r-41)/15,\,(r+4\varepsilon k-11)/5\big)$, with no condition on the number of
classes. \emph{Sketch:} solve the three bit constraints of the $Y$-step globally by Lemma (C), using four disjoint bases
found by matroid union. (496/496 constructions checked.)
\item \textbf{S0.} A shift-0 layer $C_0\cap\binomV k$ ($C_0$ even, $4\mid k$) with $\tau>3k/4$ and $n<11k/6$ has a
Fano-type seven. \emph{Sketch:} Lemma (C) twice gives a Schur triangle $x_1+x_2+x_3=0$ with empty triple intersection
and union $3k/2$; Lemma~\ref{lem:NC1} with $W=V$, $d=r$ finishes. (319/319 sevens checked.)
\item \textbf{S1.} The same with $x_1+x_2+x_3=\Omega$ gives a Fano-type seven at every shift when $r\le4\varepsilon k-7$,
at codimension up to about $k/2$ (untested).
\end{itemize}
Open for coset families: $r\ge4\varepsilon k$ with codimension above B1's bound, and $N\ge1.9385$.
\where{\file{notes/PROGRAMME2_B.md} \S3, \S8; \file{reports/RESPONSE_TO_DIRECTION_3.tex}.}

\paragraph{Bridge results \st{HAND}.}
\begin{itemize}[leftmargin=*]
\item \textbf{C1.} Three members with empty common intersection and $\wt(x_1+x_2+x_3)\le4T-3k-6$ in a $T$-covering
family give a Fano-type seven (colour by the triangle and split evenly). With a covering triple of weight $\le2r+1$ this
gives $\tau\le(2n+k)/6+2$ for every $k$-uniform $(7,2)$-family. The bound is the same as Theorem~\ref{thm:leanA}; the
new content is the robust pencil-freeness $\wt(x_1+x_2+x_3)\ge4\varepsilon k-5$ of any seven-free family.
\item \textbf{C2.} Two-block families $\{B:|B\cap X|\in J\}$ with $3k/2\le n<9k/4$ and $\tau\ge3k/4+3$ have a good
seven. (600/600 checked.)
\end{itemize}
\where{\file{notes/PROGRAMME2_C.md}.}

\section{Small values}\label{sec:small}
$f(6,7)=5$ and $f(7,7)=6$ \st{EXACT}. \emph{Lower bounds:} $K^{(k)}_n$ is $(7,2)$ iff no seven $(n-k)$-sets cover all
pairs of $[n]$, and $\tau(K^{(k)}_n)=n-k+1$. Schönheim's bound~\cite{Sch} gives $C(10,4,2)\ge\lceil\frac{10}4\cdot3\rceil=8$
and $C(12,5,2)\ge\lceil\frac{12}5\cdot3\rceil=8$, so $K^{(6)}_{10}$ and $K^{(7)}_{12}$ are $(7,2)$. \emph{Upper bounds:}
computer-generated strategies for an 8-request game at $(k,T)=(6,5),(7,6)$, replayed by an independent checker
(88 and 207 nodes), combined with the anchored-reduction lemma, which is \st{VERIFIED} for anchor-keeping strategies such as these. Also $f(8,7)\in\{6,7\}$: the
lower bound from $f(7,7)$ by adding a private point to each member, the upper bound $\lceil7\cdot8/8\rceil$ from FKW.
\where{\file{notes/verify/small_values.md}; \file{notes/FRESH_C.md} \S14.3; \file{notes/verify/TWIN13_16A.md}.}

\section{What remains}\label{sec:open}

\begin{proposition}[reduction to (S)]\label{prop:S}
(G) holds iff (S) holds. \st{consequence of VERIFIED results}
\end{proposition}
\begin{proof}
(G) implies (S) trivially. Conversely, given a family as in (G), pad it to $k$-uniform (Lemma~\ref{lem:pad}; a good
seven of the padded family restricts to one of the original). Then $\tau\ge T+1\ge(3/4+\varepsilon)k+1$. If
$\tau^*_b\le7/3-\eta$, Theorem~\ref{thm:RA} applies for $k$ large; if there is a block, Corollary~\ref{cor:NC}
applies; otherwise (S) applies.
\end{proof}

\paragraph{Open problems.}
\begin{enumerate}[leftmargin=*]
\item \textbf{(S)}, \S\ref{sec:summary}. It is a supersaturation statement at exponentially small density: covering
families have density $e^{-\gamma k}$ in the $k$-slice, with $\gamma\approx0.10$--$0.29$ in the Fano window. Minimum
coverings can be configuration-free, so minimality alone cannot give it.
\item \textbf{INV-P} (structure form, for $3\varepsilon\le\delta<1/12$). A thin covering family with no near-complete
block, in which every empty-intersection triple has $\wt(x_1+x_2+x_3)\ge4\varepsilon k-O(1)$, has all but $o(|H|)$ members
in one coset. With B this would give (S) in the Fano range. The candidate obstruction is a non-linear sum-free set of
density $\le5/16$ in a quotient with heavy triple-sum classes; none is known.
\item \textbf{BAND.} Does every $T$-covering $k$-uniform family on $(7/4+\delta)k$ points have two members with
$|x_1\cap x_2|\in[(k-r)/2,\,(k-r)/2+(5\delta/14+\varepsilon/2)k]$? A positive answer kills every Lemma-Q-type
counterexample.
\item \textbf{$N\ge11/6$.} Good sevens need not be of Fano type there, and nothing beyond Theorem~\ref{thm:leanA} is
known: LPs over all 716 seven types with at most two forced members reproduce exactly the one-anchor bound. The
obstacle is forcing a pencil, not finishing.
\end{enumerate}

\paragraph{Evidence \st{computed; small $k$ only}.} Searches that minimise the number of good sevens among covering
families at $(k,T,n)=(8,7,16),(10,8,21),(12,10,23)$ stay unstructured, with normalised seven count $1.00\pm0.02$;
planted linear layers break up. Deep cosets of BCH and Reed--Muller codes have $\tau$ far below $3k/4$. No family with
$N\ge11/6$ and $\tau>3k/4+1$ avoiding good sevens is known. This hints that (S), not a structure theorem, is the
right target. It is a hint, not a decision: no Fano-range instance is within reach.
\where{\file{notes/PROGRAMME3_TASKB.md}, \file{notes/PROGRAMME3_GAMMA.md}, \file{notes/PROGRAMME3_CX.md},
\file{notes/PROGRAMME3_BULK.md}, \file{notes/PROGRAMME2_C.md} \S6.}

\paragraph{Dead routes} (details and proofs of failure in \file{reports/FAILED_ROUTES.pdf}): bounded-request games;
junta approximation and spread lemmas (the families are less than $2.28$-spread); fractional relaxation plus rounding;
hypergraph containers with supersaturation; flag algebras and SDP on the limit object; the entropy method; covering-design
stability; linear-Fourier/BLR gluing of local affine sevens (quadratic level sets contain no Fano plane yet look
random); Fourier counting plus density increment for sparse families (the $\gamma$ wall).

\begin{thebibliography}{9}\small\raggedright
\bibitem{EP} Erd\H{o}s Problems website, Problem \#644, \url{https://www.erdosproblems.com/644}.
\bibitem{EFKT} P.~Erd\H{o}s, D.~Fon-Der-Flaass, A.~V.~Kostochka, Zs.~Tuza, Small transversals in uniform
hypergraphs, \emph{Siberian Adv.\ Math.}\ 2 (1992) 82--88. [cite?: not seen; values as reported in~\cite{EP}]
\bibitem{FKW} D.~G.~Fon-Der-Flaass, A.~V.~Kostochka, D.~R.~Woodall, Transversals in uniform hypergraphs with property
$(7,2)$, \emph{Discrete Math.}\ 207 (1999) 277--284.
\bibitem{NP} D.~T.~Nagy, B.~Patk\'os, Triangles in intersecting families, arXiv:2201.02452.
\bibitem{BDMS} A.~Bishnoi, S.~Das, P.~Morris, T.~Szab\'o, Ryser's conjecture for $t$-intersecting hypergraphs,
arXiv:2001.04132.
\bibitem{OEIS} OEIS Foundation, Sequence A001206 (maximal intersecting families), \url{https://oeis.org/A001206}.
\bibitem{Sch} J.~Sch\"onheim, On coverings, \emph{Pacific J.\ Math.}\ 14 (1964) 1405--1411. [cite?]
\end{thebibliography}
\end{document}
